\section{Results}
\label{sec:results}

\subsection{Complete escalations}

Over the \nChainDenominator units the chain analysis covered, it found
\nChainsFound complete cross-agent escalations. With no event observed in
\nChainDenominator trials, the one sided 95\% upper bound on the rate is
\chainBound\%. The measurement is consistent with a true rate anywhere from
zero to about one deployment in ten.

The predicate does fire inside the wider set of \nProjects analysed projects.
\nEscalating configurations contain an escalation, \nEscalatingOwnResource of
them reaching a resource the same configuration itself writes, and
\nSurvivingTakeover of them retaining a takeover edge.

\nEscalatingSingleAgent of those \nEscalating configurations declare a single
agent, and the escalating principals are a continuous integration
deploy role, a service user, and a role that executes actions. They are single-identity
escalations, produced by the same predicate.

\subsection{Candidates blocked on one permission}

\nBlocked takeover candidates across \nBlockedConfigs configurations are
blocked, and every one of them is blocked by the same permission,
\iam{iam:PassRole}. This is the structural claim of
Section~\ref{sec:background} appearing in the measurement: there is no second
way around the pass.

The \nBlocked are not interchangeable. Figure~\ref{fig:blocked} sorts them by
what the principal and the target are, recovered from the raw Terraform rather
than from the analysis output. \nBlockedAgentToAgent are one agent reaching
another agent of the same deployment, which is the case this paper is about.
\nBlockedSingleAgent are in configurations that declare one agent, so the
precondition is present and the second agent is not. \nBlockedPrincipalNotAgent
have a principal that is not an agent execution role, such as a build role, and
\nBlockedTargetNotAgent have a target that is not one.

\begin{figure}[t]
\centering
\resizebox{\columnwidth}{!}{%
\begin{tikzpicture}
\begin{axis}[
  width=72mm,
  height=32mm,
  xbar stacked,
  xmin=0, xmax=26,
  bar width=7mm,
  axis y line=none,
  axis x line*=bottom,
  xtick={0,5,10,15,20,25},
  xlabel={\small blocked takeover candidates},
  xlabel style={font=\small, yshift=1mm},
  tick label style={font=\scriptsize},
  legend style={
    font=\scriptsize, draw=none, fill=none,
    at={(0.5,1.28)}, anchor=north, legend columns=2,
    /tikz/every even column/.append style={column sep=3mm}},
  nodes near coords={\scriptsize\pgfmathprintnumber\pgfplotspointmeta},
  every node near coord/.append style={text=white, anchor=center},
]
\addplot+[fill=cStored, draw=cStored] coordinates {(10,0)};
\addlegendentry{one agent reaching another}
\addplot+[fill=cNeutral, draw=cNeutral] coordinates {(10,0)};
\addlegendentry{configuration declares one agent}
\addplot+[fill=cIdentity, draw=cIdentity] coordinates {(4,0)};
\addlegendentry{principal is not an agent role}
\addplot+[fill=cAccent, draw=cAccent] coordinates {(2,0)};
\addlegendentry{target is not an agent role}
\end{axis}
\end{tikzpicture}}
\caption{The \nBlocked candidates blocked on \iam{iam:\allowbreak PassRole}, sorted by what
the principal and the target are. Only the first group,
\nBlockedAgentToAgent of \nBlocked, is one agent reaching another agent of the
same deployment. Each candidate was re-derived from the raw Terraform rather
than taken from the analysis output.}
\label{fig:blocked}
\end{figure}


Separately, \nPassRoleWide configuration grants \iam{iam:PassRole} broadly
enough to cover a role an agent in that same configuration runs as. That is the
one deployment in the corpus where the gate that blocks the other
\nBlocked candidates is already open.

\subsection{What an existing scanner reports on the same corpus}
\label{sec:baseline}

The property here is a relation between two resources, and rule based
scanners work on one. We ran Checkov, a widely used open source scanner for
infrastructure as code, over the same \nProjects configurations, on Terraform
with no network access and nothing else changed.

Checkov reports \nCheckovFailed findings from \nCheckovRules distinct rules,
and \nCheckovOnAgentRuntime of them take an agent runtime as their subject:
the resource type is not in its rule set, so these agents are invisible to it.
Its rules that span resources fired \nCheckovGraphFailed times.
\nCheckovRelationRules of the rule names that fired suggest a relation
between resources, and they concern cross-account access, cross-region
replication and load balancing, not two agents.

Most of those findings are about encryption and logging, which this paper
does not look at. A deployment can pass every rule in a mature rule set while
\nSharedRoleRuntimes of \nMultiAgentRuntimes multi-agent deployments run every
agent under one role.

\subsection{Agents sharing one execution role}

\begin{table}[t]
\centering
\caption{Multi-agent deployments and role sharing, counted two ways. A runtime
block carrying \texttt{count} or \texttt{for\_each} deploys more than one
agent, so the two columns answer different questions. The paper uses the
runtime column, and the artifact names the deployments that move.}
\label{tab:sharing}
\small
\setlength{\tabcolsep}{4pt}
\begin{tabular}{@{}lcc@{}}
\toprule
 & \textbf{by block} & \textbf{by runtime} \\
\midrule
Two or more agents               & \nMultiAgentBlocks & \nMultiAgentRuntimes \\
Of those, all on one role        & \nSharedRoleBlocks & \nSharedRoleRuntimes \\
\bottomrule
\end{tabular}
\end{table}

Table~\ref{tab:sharing} gives the role sharing census. \nSharedRoleRuntimes of
\nMultiAgentRuntimes multi-agent deployments run every agent under a single
execution role.

This is the finding with the most practical weight in the paper, and it is not
an escalation. Where the role is shared the attacker does not need the takeover
at all: the agents already hold identical permissions, so compromising any one
of them yields what all of them can reach. It also removes the account's
ability to attribute a call to an agent, because every call arrives under the
same role.

\subsection{CloudFormation}

The CloudFormation side of the collection holds \nCfnMultiRuntime templates
that declare more than one runtime. \nCfnAnalysed were analysed and contain
\nChainsFound chains. \nCfnUndecidable are undecidable: an intrinsic in a role
or runtime argument did not resolve, and the template is reported as such
rather than as clean.

\subsection{CDK}

The source scan flags \nCdkMultiRuntimeDirs CDK directories with more than one
runtime, \nCdkInScope after the exclusion rule of
Section~\ref{sec:methodology}. Reading a CDK role's policies out of the source
is a heuristic, so we synthesised instead.

Of the \nCdkInScope, \nCdkAnalysed synthesised and went through the exact
CloudFormation pipeline, recovering \nCdkRuntimes agent runtimes and
\nCdkEscalating escalations. Table~\ref{tab:cdk} gives the reasons the rest did
not synthesise. A failure to install a dependency and a failure to find an
application are different facts about the corpus.

\begin{table}[t]
\centering
\caption{Why \nCdkInScope CDK directories did or did not reach the analysis.
A directory that does not synthesise is not counted as clean.}
\label{tab:cdk}
\small
\begin{tabular}{lr}
\toprule
\textbf{Outcome} & \textbf{Directories} \\
\midrule
Synthesised and analysed                     & \nCdkAnalysed \\
\midrule
Synthesis failed                             & \nCdkFailSynth \\
Not a supported application layout           & \nCdkFailLanguage \\
Dependency install failed                    & \nCdkFailInstall \\
No \texttt{cdk.json} above the directory     & \nCdkFailLocate \\
\bottomrule
\end{tabular}
\end{table}

\subsection{What the CDK results add to the denominator}

\nCdkAnalysed directories is not \nCdkAnalysed applications.
\nCdkSampleCopies of them are the same AWS sample application carried into
\nCdkSampleCopies different repositories, at the same path inside each. By the
deduplication rule of Section~\ref{sec:methodology} those count once, which
leaves \nCdkDistinctApps distinct applications.

Adding them to the chain denominator gives \nChainDenominatorExt units with
\nChainsFound chains, and the one sided 95\% upper bound tightens from
\chainBound\% to \chainBoundExt\%. The first bound covers the deployments read
directly; the second adds the synthesised CDK applications.
